\chapter{Introduction}

\section{The thesis}

\textbf{Hardware is asynchronous; the operating system and the userspace
runtime should be too.} Two projects, started independently from opposite ends of
the software stack, built toward the same shape and converged on the same
primitive: a waitable completion handle.

\paragraph{CharlotteOS} starts in the kernel and asks: \emph{hardware is
asynchronous with respect to the CPU, so why is the OS synchronous?} Its answer
is to make the kernel async-first --- cheap threads, an observer/waker event
model, and system calls that are asynchronous by default, returning a
completion capability that userspace can wait on.

\paragraph{sitas} starts in userspace and asks: \emph{what does shared-nothing,
shard-per-core concurrency look like when it is built out of Rust ownership
instead of locks?} Its answer is a discipline: each shard owns its state, only
the owning shard mutates it, and everything that crosses a shard boundary is an
owned value moved through an explicit typed message.

Both projects independently converged. sitas's \texttt{Reply<T>} and
\texttt{Notify} (awaitable completions) are the same object as CharlotteOS's
\texttt{Observer} and \texttt{Waker}, just with different names. sitas's shard
mailbox (\texttt{ShardSender}/\texttt{ShardReceiver}) is a bounded per-core
typed queue --- as are CharlotteOS's bounded per-LP \texttt{IPI\_CMD\_QUEUES}.
sitas's \texttt{ShardLocal<T>} (lock-free, owner-checked, closure-based access)
is a Rust-ownership alternative to CharlotteOS's \texttt{PerLp<T>}
(lock-wrapped).

This document describes what happens when these two projects meet: a userspace
runtime whose async model \textbf{agrees} with the kernel it runs on, instead
of emulating async-first on a blocking foundation. That agreement removes the
``retrofit tax'' --- the async-signal-safety, thread pools, and
\texttt{io\_uring}-over-blocking emulation that exist only because the runtime
and the OS disagree about what is fundamental.

\section{What was built and verified}

The collaboration produced a booting, self-testing async-first kernel on both
AArch64 and x86-64, with a full async syscall ABI (completion capabilities, CQ
ring, exit-observer completion), a sitas-derived shard-per-core programming
model in the kernel (\texttt{ShardLocal<T>}, \texttt{ShardMailbox<M>}, bounded
IPI queues), and a sitas shard-per-core executor running as a user binary on
CharlotteOS. Everything was verified on QEMU.

\textbf{Key deliverables:}

\begin{itemize}
    \item \textbf{Completion-capability subsystem} --- per-address-space
        capability tables, submit / complete / poll / wait / cancel / close,
        \texttt{submit\_worker} (exit-observer), shared-memory CQ ring for
        zero-syscall completion delivery.
    \item \textbf{Syscall dispatch} --- AArch64 SVC handler decodes
        \texttt{ESR\_EL1.EC}, builds a \texttt{TrapFrame}, dispatches to a
        typed handler table. x86-64 SYSCALL trampoline (MSR-based entry,
        \texttt{swapgs} / \texttt{sysretq}).
    \item \textbf{Scheduler} --- \texttt{RoundRobin} per LP, quantum timer,
        block / yield / wake, deferred thread reaper, per-LP thread pinning
        (\texttt{submit\_to\_lp} for sitas shard placement).
    \item \textbf{Cross-LP IPI} --- bounded per-LP queues with first-class
        backpressure, typed closure dispatch \\ (\texttt{IpiRpc::Closure}).
    \item \textbf{Shard-per-core programming model} ---
        \texttt{ShardLocal<T>} (lock-free, owner-checked), \\
        \texttt{ShardMailbox<M>} (typed bounded queue).
    \item \textbf{Real-EL0 user thread} --- compiled Rust binary runs at EL0,
        calls SVC from userspace, round-trips through the kernel's dispatch
        and returns.
    \item \textbf{sitas integration} --- sitas executor (no\_std) +
        sitas-charlotte ReactorBackend + \\ \texttt{ShardRuntime::spawn\_shard}
        via SVC \#7. The sitas shard executor runs on CharlotteOS.
\end{itemize}

\section{How this document is organized}

\begin{itemize}
    \item \textbf{Chapter 2} --- Architecture: Logical Processors, Observers,
        capability tables, the scheduler, IPIs.
    \item \textbf{Chapter 3} --- Execution model: threads, blocking, context
        switching, timers.
    \item \textbf{Chapter 4} --- The async syscall ABI: the five operations,
        exit-observer completions, the CQ ring, cancellation, backpressure.
    \item \textbf{Chapter 5} --- Programming model: invariants, ShardLocal,
        ShardMailbox, guidance.
    \item \textbf{Chapter 6} --- Development: build, QEMU boot, self-tests, CI.
    \item \textbf{Chapter 7} --- What we learned: bugs found on hardware, root
        causes, fixes, and key insights.
    \item \textbf{Chapter 8} --- Current status: what is verified, known
        limitations, repositories.
\end{itemize}

\section{Who this is for}

Developers familiar with Rust (ownership, \texttt{Send}, \texttt{Future}),
concurrency primitives, and basic OS architecture. Every claim is based on what
was built and verified on QEMU AArch64 and x86-64. Unimplemented ideas are
flagged.

\endinput
